% !TEX program = lualatex
% ============================================================
%  Stat 150 — Comprehensive Solutions
%  Primary: Durrett, Essentials of Stochastic Processes (3rd ed.)
%  Secondary: Pinsky & Karlin, An Introduction to
%             Stochastic Modeling (4th ed.)
%  Keshav Ramamurthy — Fall 2026
% ------------------------------------------------------------
%  HOW TO USE (one growing document, one banner per chapter,
%  two parts: Part I Durrett, Part II Pinsky--Karlin)
%
%    Find the chapter banner below and type directly under it:
%
%        \exercise{1.4.2}
%        \begin{solution}
%          your write-up ...
%        \end{solution}
%
%    Restate a problem first (optional, renders a gold box):
%
%        \begin{problem}[short title]
%          statement ...
%        \end{problem}
%
%    Structure inside a solution:  claim / idea / proofsketch
%    environments, plus defn, thm, lem, prop, cor, ex, rem
%    (numbered within the current chapter).
%    Cross-reference:   \exref{1.4.2}
% ============================================================
\documentclass[11pt]{article}

\usepackage{fontspec}
\usepackage{microtype}
\usepackage{mathtools,amssymb,amsfonts,amsthm,bm}
\usepackage{enumitem}
\usepackage{booktabs,array,xcolor}
\usepackage{geometry,fancyhdr,setspace}
\usepackage{etoolbox}
\usepackage[most]{tcolorbox}
\usepackage[hidelinks]{hyperref}
\usepackage[nameinlink,noabbrev]{cleveref}

% ---------- Page and paragraph rhythm (matches lecture template) ----------
\geometry{margin=1in,headheight=15pt}
\setstretch{1.12}
\setlength{\parindent}{0pt}
\setlength{\parskip}{0.55em plus 0.12em minus 0.08em}
\setlength{\abovedisplayskip}{0.8em plus 0.2em minus 0.1em}
\setlength{\belowdisplayskip}{0.8em plus 0.2em minus 0.1em}
\allowdisplaybreaks[3]
\setlist{leftmargin=*,topsep=0.35em,itemsep=0.15em,parsep=0pt}

% ---------- Palette (matches reference cheatsheet) ----------
\definecolor{ink}{HTML}{1E293B}
\definecolor{navy}{HTML}{183B56}
\definecolor{paperblue}{HTML}{E8F1F8}
\definecolor{gold}{HTML}{FFF3CD}
\definecolor{mint}{HTML}{E8F5E9}
\definecolor{muted}{HTML}{52606D}
\definecolor{rule}{HTML}{C9D5DF}
\color{ink}

% ---------- Metadata ----------
\newcommand{\studentname}{Keshav Ramamurthy}
\newcommand{\coursename}{Stat 150}
\newcommand{\courselongname}{Stochastic Processes}
\newcommand{\textbookname}{Durrett, \emph{Essentials of Stochastic Processes} (3rd ed.)}

% ---------- Core notation (same as ~/.config/nvim/textemplate.tex) ----------
\newcommand{\N}{\mathbb{N}}
\newcommand{\Z}{\mathbb{Z}}
\newcommand{\Q}{\mathbb{Q}}
\newcommand{\R}{\mathbb{R}}
\newcommand{\C}{\mathbb{C}}
\newcommand{\F}{\mathbb{F}}
\newcommand{\K}{\mathbb{K}}
\newcommand{\E}{\mathbb{E}}
\newcommand{\Prob}{\mathbb{P}}
\newcommand{\1}{\mathbf{1}}
\newcommand{\eps}{\varepsilon}
\newcommand{\dd}{\,\mathrm{d}}
\newcommand{\abs}[1]{\left\lvert#1\right\rvert}
\newcommand{\norm}[1]{\left\lVert#1\right\rVert}
\newcommand{\ip}[2]{\left\langle#1,#2\right\rangle}
\newcommand{\set}[1]{\left\{#1\right\}}
\newcommand{\ceil}[1]{\left\lceil#1\right\rceil}
\newcommand{\floor}[1]{\left\lfloor#1\right\rfloor}
\newcommand{\given}{\,\middle\vert\,}
\newcommand{\indep}{\perp\!\!\!\perp}
\newcommand{\wh}[1]{\widehat{#1}}
\newcommand{\deriv}[2]{\frac{\mathrm{d}#1}{\mathrm{d}#2}}
\newcommand{\pderiv}[2]{\frac{\partial#1}{\partial#2}}
\newcommand{\lto}{\longrightarrow}
\newcommand{\st}{\text{ such that }}
\DeclareMathOperator{\Aut}{Aut}
\DeclareMathOperator{\End}{End}
\DeclareMathOperator{\Gal}{Gal}
\DeclareMathOperator{\Hom}{Hom}
\DeclareMathOperator{\Ker}{ker}
\DeclareMathOperator{\im}{im}
\DeclareMathOperator{\rank}{rank}
\DeclareMathOperator{\nullity}{nullity}
\DeclareMathOperator{\Span}{span}
\DeclareMathOperator{\tr}{tr}
\DeclareMathOperator{\diag}{diag}
\DeclareMathOperator{\spec}{spec}
\DeclareMathOperator{\ord}{ord}
\DeclareMathOperator{\supp}{supp}
\DeclareMathOperator{\diam}{diam}
\DeclareMathOperator{\Var}{Var}
\DeclareMathOperator{\Cov}{Cov}
\DeclareMathOperator*{\argmin}{arg\,min}
\DeclareMathOperator*{\argmax}{arg\,max}

% ---------- Course-specific notation (stochastic processes) ----------
\DeclareMathOperator{\Expo}{Exp}                    % exponential distribution
\DeclareMathOperator{\Poi}{Poisson}                % Poisson distribution
\DeclareMathOperator{\Bin}{Binomial}
\DeclareMathOperator{\Geom}{Geometric}
\newcommand{\Ex}[1]{\E\!\left[#1\right]}           % E[ ... ]
\newcommand{\Prb}[1]{\Prob\!\left(#1\right)}       % P( ... )
\newcommand{\trans}[1]{p_{#1}}                     % transition probability p_ij
\newcommand{\rate}[1]{\lambda_{#1}}                % generator rate

% ---------- Theorem-like environments (numbered by chapter) ----------
\newcounter{chap}
\newtheoremstyle{notestyle}{0.7em}{0.7em}{\itshape}{}{}{\bfseries}{.5em}{}
\theoremstyle{notestyle}
\newtheorem{theorem}{Theorem}[chap]
\newtheorem{lemma}[theorem]{Lemma}
\newtheorem{proposition}[theorem]{Proposition}
\newtheorem{corollary}[theorem]{Corollary}
\theoremstyle{definition}
\newtheorem{definition}[theorem]{Definition}
\newtheorem{example}[theorem]{Example}
\newtheorem{remark}[theorem]{Remark}

% ---------- Problem machinery ----------
% Chapter banner: \chapterbanner{1}{Markov Chains}{Durrett, \S 1.1--1.13}
\newtcolorbox{chapterbox}[2]{%
  enhanced,
  colback=paperblue, colframe=navy, boxrule=0.7pt, arc=2.5pt,
  left=10pt, right=10pt, top=12pt, bottom=8pt,
  boxed title style={colback=navy, colframe=navy, arc=2pt, boxrule=0pt,
    left=8pt, right=8pt, top=3pt, bottom=3pt},
  attach boxed title to top left={yshift=-2.4mm, xshift=5mm},
  title={\sffamily\bfseries\color{white}\large Chapter #1\quad #2}}
\newcommand{\chapterbanner}[3]{%
  \clearpage
  \setcounter{chap}{#1}%
  \setcounter{theorem}{0}%
  \phantomsection
  \addcontentsline{toc}{section}{Chapter #1\quad #2}%
  \begin{chapterbox}{#1}{#2}
    {\small\sffamily\color{muted}#3\par}
  \end{chapterbox}
  \medskip\nopagebreak}

% Exercise header: \exercise{1.4.2}  →  bold header + thin rule, then type
\newcommand{\exercise}[1]{%
  \par\goodbreak\medskip\noindent
  {\sffamily\bfseries\color{navy}Exercise #1}%
  \par\nobreak\vspace{1.5pt}
  {\color{rule}\hrule height 0.7pt}\par\nobreak\smallskip\nopagebreak}

% Cross-reference an exercise: \exref{1.4.2}
\newcommand{\exref}[1]{\textsf{\textbf{Exercise~#1}}}

% Problem statement box: \begin{problem}[short title] ... \end{problem}
\newtcolorbox{problem}[1][]{%
  enhanced, breakable,
  colback=gold!50, colframe=rule, boxrule=0.5pt, arc=2pt,
  colbacktitle=gold!50, coltitle=navy,
  left=8pt, right=8pt, top=6pt, bottom=6pt,
  before skip=8pt, after skip=8pt,
  title={\sffamily\bfseries Problem\ifstrempty{#1}{}{\ ---\ #1}}}

% Solution / answer blocks (qed square at the end)
\newenvironment{solution}{\begin{proof}[Solution]}{\end{proof}}
\newenvironment{answer}{\begin{proof}[Answer]}{\end{proof}}
\renewcommand{\qedsymbol}{$\square$}

% Quick structure inside solutions
\newenvironment{claim}{\par\smallskip\noindent\textbf{Claim.}\ }{\par\smallskip}
\newenvironment{idea}{\par\smallskip\noindent\textit{Idea.}\ }{\par\smallskip}
\newenvironment{proofsketch}{\par\smallskip\noindent\textbf{Proof sketch.}\ }{\par\smallskip}

% ---------- Header ----------
\pagestyle{fancy}
\fancyhf{}
\fancyhead[L]{\coursename}
\fancyhead[C]{\courselongname\ — solutions}
\fancyhead[R]{\studentname}
\fancyfoot[C]{\thepage}

\begin{document}

% ---------- Title ----------
\begin{center}
  {\Large\sffamily\bfseries\color{navy}\coursename\ -- \courselongname}\\[3pt]
  {\sffamily\color{muted} Comprehensive exercise solutions}\\[3pt]
  {\small\sffamily\color{muted}\textbookname\quad\textbullet\quad Fall 2026\quad\textbullet\quad \studentname}
\end{center}
\medskip
{\color{rule}\hrule height 0.8pt}
\bigskip

{\small
\begin{tcolorbox}[colback=paperblue, colframe=rule, boxrule=0.4pt, arc=2pt,
  left=7pt, right=7pt, top=5pt, bottom=5pt,
  title={\sffamily\bfseries\color{navy} Adding a solution},
  colbacktitle=paperblue]
\sffamily To record a solved exercise, type under its chapter banner:
\begin{center}
\texttt{\textbackslash exercise\{1.4.2\}\ \ \textbackslash begin\{solution\}\ ...\ \textbackslash end\{solution\}}
\end{center}
Restate the problem first with \texttt{\textbackslash begin\{problem\}[title]...\textbackslash end\{problem\}};
structure proofs with \texttt{claim}, \texttt{idea}, \texttt{proofsketch}.
\end{tcolorbox}
}

\medskip
\tableofcontents

% =====================================================================
%  PART I — DURRETT, ESSENTIALS OF STOCHASTIC PROCESSES (3rd ed.)
% =====================================================================
\part*{Part I\quad Durrett, \emph{Essentials of Stochastic Processes}}
\addcontentsline{toc}{part}{Part I — Durrett, Essentials of Stochastic Processes}

% ---------------------------------------------------------------------
\chapterbanner{1}{Markov Chains}{Durrett, \S 1.1--1.13:\ definitions, multistep transitions, classification of states, stationary distributions, limit behavior, exit distributions and times}

% \exercise{1.1.1}
% \begin{solution}
%   ...
% \end{solution}

% ---------------------------------------------------------------------
\chapterbanner{2}{Poisson Processes}{Durrett, \S 2.1--2.6:\ exponential distribution, constructing the process, compound Poisson, thinning/superposition/conditioning}

% ---------------------------------------------------------------------
\chapterbanner{3}{Renewal Processes}{Durrett, \S 3.1--3.5:\ laws of large numbers, queueing applications, age and residual life}

% ---------------------------------------------------------------------
\chapterbanner{4}{Continuous Time Markov Chains}{Durrett, \S 4.1--4.8:\ generators, transition probabilities, limiting behavior, exit times, Markovian queues}

% ---------------------------------------------------------------------
\chapterbanner{5}{Martingales}{Durrett, \S 5.1--5.5:\ conditional expectation, gambling strategies and stopping times, applications}

% ---------------------------------------------------------------------
\chapterbanner{6}{Mathematical Finance}{Durrett, \S 6.1--6.3:\ binomial model, Black--Scholes, options}

% =====================================================================
%  PART II — PINSKY & KARLIN, AN INTRODUCTION TO STOCHASTIC MODELING
% =====================================================================
\part*{Part II\quad Pinsky \& Karlin, \emph{An Introduction to Stochastic Modeling}}
\addcontentsline{toc}{part}{Part II — Pinsky \& Karlin, An Introduction to Stochastic Modeling}

% ---------------------------------------------------------------------
\chapterbanner{1}{Introduction}{Pinsky \& Karlin:\ stochastic modeling, probability review}

% ---------------------------------------------------------------------
\chapterbanner{2}{Conditional Expectation and Related Topics}{Pinsky \& Karlin:\ conditioning, generating functions, martingales}

% ---------------------------------------------------------------------
\chapterbanner{3}{Markov Chains}{Pinsky \& Karlin:\ discrete-time chains, first-step analysis}

% ---------------------------------------------------------------------
\chapterbanner{4}{The Long Run Behavior of Markov Chains}{Pinsky \& Karlin:\ stationary and limiting distributions, absorption}

% ---------------------------------------------------------------------
\chapterbanner{5}{Poisson Processes}{Pinsky \& Karlin:\ the Poisson process and its transformations}

% ---------------------------------------------------------------------
\chapterbanner{6}{Continuous-Time Markov Chains}{Pinsky \& Karlin:\ birth--death processes, generators, branching processes}

% ---------------------------------------------------------------------
\chapterbanner{7}{Renewal Processes}{Pinsky \& Karlin:\ renewal theory and applications}

% ---------------------------------------------------------------------
\chapterbanner{8}{Brownian Motion and Gaussian Processes}{Pinsky \& Karlin:\ Brownian motion, drift, variants}

% ---------------------------------------------------------------------
\chapterbanner{9}{Queueing Systems}{Pinsky \& Karlin:\ Markovian queues, Little's formula, networks}

\end{document}
